% DEFINED:   sec:bengio-arch, sec:embedding-layer
% RESOLVED:  ch:foundations, sec:mlp-backprop (backward into Part I)
% FORWARD:   (none)

\section{The architecture}\label{sec:bengio-arch}

The model has three parts, and each one is a building block already in
hand or close to it.

First, an \textbf{embedding table} $E$ of shape $V \times d$, where $V$ is
the vocabulary size and $d$ is the embedding dimension. Each token id $c$
indexes a row $E_c$, a learned dense vector of length $d$. The rows are
\emph{shared across positions}: the embedding of a token is the same
vector whether the token appears first or last in the context window.

Second, a \textbf{concatenation step}. Given a context of $N$ token ids
$c_{t-N+1}, \ldots, c_t$, look up the $N$ embeddings and concatenate them
into a single vector $\vx$ of length $Nd$. This vector carries two kinds
of information at once: \emph{what} tokens are present, through the
embeddings, and \emph{where} each one sits in the window, through its
offset in the concatenation.

Third, a \textbf{small multilayer perceptron} that maps $\vx$ to logits
over the vocabulary,
\[
  \mathbf{h} = \tanh(W_1 \vx + \mathbf{b}_1), \qquad
  \vz = W_2 \mathbf{h} + \mathbf{b}_2,
\]
exactly the head from \Cref{ch:foundations}. A softmax of $\vz$ is the
predicted distribution over the next token, and
\texttt{SoftmaxCrossEntropy} is the loss. The network emits logits; the
loss interprets them, the same division of labor the whole book has kept.

That is the entire model. In \texttt{scratchnn} it is a plain
\texttt{Network} with no custom training loop:
\begin{lstlisting}
net = Network([
    EmbedConcat(vocab_size, embed_dim, context_size),
    Linear(context_size * embed_dim, hidden_size),
    Tanh(),
    Linear(hidden_size, vocab_size),
], SoftmaxCrossEntropy())
\end{lstlisting}
The only new layer is \texttt{EmbedConcat}, which wraps the embedding
lookup and the concatenation. Its forward takes a \texttt{list[int]} of
length \texttt{context\_size} and returns a flat \texttt{list[float]} of
length \texttt{context\_size * embed\_dim}. From the downstream layers'
point of view the input is just a vector, so \texttt{Linear},
\texttt{Tanh}, and \texttt{SoftmaxCrossEntropy} compose without change.
Training is \texttt{net.fit(X, Y, ...)} where each \texttt{X[i]} is a
context (a \texttt{list[int]}) and each \texttt{Y[i]} is the next token id.

\section{The Embedding layer}\label{sec:embedding-layer}

An \texttt{Embedding} maps a discrete token id to a dense vector. It is
mathematically a \texttt{Linear} layer applied to a one-hot input, but it
skips the multiply-by-zero by indexing the row directly. The forward
returns a copy of the relevant row and records the id; the backward pops
the id and accumulates the upstream gradient into that row. Here is the
real layer from \texttt{scratchnn}, with the constructor trimmed:
\begin{lstlisting}
class Embedding(Layer):
    def forward(self, token_id):
        self.cache.append(token_id)
        return list(self.weights[token_id])

    def backward(self, grad):
        token_id = self.cache.pop()
        row = self.dweights[token_id]
        for k in range(len(grad)):
            row[k] += grad[k]
        return None    # discrete input, no upstream gradient
\end{lstlisting}
The lookup view and the one-hot view are the same computation. Real
frameworks (PyTorch's \texttt{nn.Embedding}, for one) implement the lookup
form for efficiency; the math does not change. The weight table is exactly
the matrix a \texttt{Linear} layer would hold for one-hot inputs.

Because the same \texttt{Embedding} is called once per context position in
a single forward pass, the layer keeps an internal LIFO cache of which ids
were used. The \texttt{backward} pops the cache and accumulates each
gradient into the matching row, the same \texttt{+=} accumulation pattern
\texttt{Linear} uses in \Cref{sec:mlp-backprop}, here routed by token id
rather than by neuron.

\texttt{EmbedConcat} is the thin wrapper that turns $N$ token ids into the
one flat vector the MLP head wants. Its forward embeds each id and
concatenates; its backward slices the incoming gradient into $N$ chunks of
length $d$ and routes each back in reverse, so the pops line up with the
pushes:
\begin{lstlisting}
class EmbedConcat(Layer):
    def forward(self, ids):           # ids: list[int] of length N
        self.embed.reset_cache()
        out = []
        for token_id in ids:
            out.extend(self.embed.forward(token_id))
        return out                    # flat list, length N * d

    def backward(self, g):
        d = self.embed_dim
        # Backward in reverse to match the embedding's LIFO cache.
        for i in reversed(range(self.context_len)):
            self.embed.backward(g[i * d:(i + 1) * d])
        return None
\end{lstlisting}
The concatenation is the whole positional prior in one line of indexing:
token $i$'s embedding always lands in slots $[id, (i+1)d)$, so the MLP head
downstream can learn position-specific weights over those slots. The
\texttt{parameters()} method delegates to the wrapped \texttt{Embedding},
so the table is the only place the rows live, and the generic optimizer
from \Cref{ch:foundations} steps them without knowing an embedding exists.
